\subsection{Clase débil intrínseca y diagonales de una sola casilla}\label{sec:B5T01}
Un residual conserva las coordenadas originales del tablero ordinario de reinas.
Su soporte $E$ tiene $k$ filas y $k$ columnas; cada casilla pertenece a
cuatro líneas y dos líneas distintas se cortan en a lo sumo una casilla.
Sea $p$ un punto estacionario interior de Bethe, con masas RC iguales a uno.
Para constantes fijas $C,\gamma,\eta$, definimos la clase determinista
\[
 p_e\le C/k,\quad \operatorname{gap}_{RC}(p)=\lambda_2(\mathscr L_p)\ge\gamma,
 \quad h_d=1-m_d\ge\eta,\quad
 m_d=\sum_{e\in d}p_e\le\varepsilon\le\varepsilon_0.
\]
El gauge RC es $(\mathbf1_R,-\mathbf1_C)$, y la brecha se toma fuera de él. Se exige además la
conectividad RC. Para $k\ge k_0$ se tiene $p<1/4$. A diferencia de la clase macroscópica, no se
exige $k\ge\delta_0n$, y las constantes no dependen del número de diagonales. La cota de $p$ obliga
a que cada fila y cada columna tengan al menos $k/C$ casillas. Una densidad global de máscara,
por sí sola, no sustituye esas condiciones.

La geometría original sigue siendo parte del problema: que la clase sea intrínseca no significa
olvidar cuáles casillas comparten diagonal.

\paragraph{Cota inferior discreta de la masa diagonal.}
Como cada fila suma uno y tiene a lo sumo $k$ casillas, resulta
\[
 \max_e p_e\ge1/k,\qquad \max_d m_d\ge1/k.
\]
En particular, una cota de la forma $\max_dm_d\le Ak/n$ exigiría $k^2\ge n/A$.

\paragraph{Eliminación exacta de una diagonal con una sola casilla.}
Escribamos $q=1-p$ y permitamos una multiplicidad $d\in\{2,3,4\}$:
\[
 \tau_d=\frac{pq}{1-dp},\qquad b_d=\frac{1-dp}{q^d}.
\]
La línea con una sola casilla tiene vacancia $q$. Su pivote y su complemento de Schur son
\[
 q+\tau_d=q\frac{1-(d-1)p}{1-dp},\quad
 \tau_d-\frac{\tau_d^2}{q+\tau_d}=\tau_{d-1},\quad
 b_d(q+\tau_d)=b_{d-1}.
\]
Por ello el producto del factor local por el determinante se conserva exactamente.
La entropía de Bernoulli de ese factor cancela el cambio de multiplicidad
en la entropía de Bethe. El gauge RC no cambia y tiene coordenada cero
en la línea eliminada. Dos líneas de una sola casilla llevan de $d=4$ a $d=2$.
Una línea sin casillas aporta un pivote igual a uno.

No se puede borrar la fila matricial manteniendo $\tau_4,b_4$. Para evitar una nueva contabilidad
de multiplicidades variables, las Secciones~\ref{sec:B5T02}--\ref{sec:B5T05} mantienen las cuatro
incidencias de cada casilla, incluso con líneas de una sola casilla. El número de estas líneas no
perjudica las cotas por masa. El caso de torres se puede representar añadiendo a cada casilla dos
diagonales de una sola casilla y eliminándolas con esta identidad.

\paragraph{Alcance.}
El predictor y el cofactor RC son los de las Secciones~\ref{sec:B3T01} y~\ref{sec:B4T01}.
No se afirma la pertenencia típica a esta clase ni la completabilidad de sus miembros.
